\chapter{PHƯƠNG PHÁP}
\label{chap:method}

Chương này mô tả toàn bộ pipeline theo thứ tự xử lý: tiền xử lý ảnh, phân tích hệ màu của biển, định vị biển số bằng đặc trưng cạnh và hình thái học, chấm điểm lại ứng viên bằng HOG và SVM tuyến tính, nắn biển về hình chữ nhật chuẩn bằng phép biến đổi phối cảnh, nhị phân hóa và tách ký tự. Dữ liệu, cách chia tập và các độ đo được trình bày ở Chương~\ref{chap:intro}; kết quả thí nghiệm ở Chương~\ref{chap:results}.

\section{Luồng xử lý}
\label{sec:flow}

Dự án chỉ dùng xử lý ảnh cổ điển, không dùng học sâu. Luồng xử lý một ảnh được tóm tắt ở Hình~\ref{fig:pipeline}. Phương pháp \code{edge} dùng trực tiếp ứng viên có điểm heuristic cao nhất; phương pháp \code{edge\_svm} xếp lại tối đa 5 ứng viên theo điểm của SVM rồi lấy ứng viên đầu. Biển được nắn về ảnh $220\times 70$, sau đó nhị phân hóa và tách ký tự.

\begin{figure}[htbp]
\centering
\begin{tikzpicture}[
  node distance=5mm and 7mm,
  box/.style={draw, rounded corners=2pt, align=center, minimum height=11mm, text width=27mm, font=\fontsize{7pt}{8pt}\selectfont},
  arr/.style={-{Stealth[length=2mm]}, thick}]
  \node[box] (in) {Ảnh xe (BGR)};
  \node[box, right=of in] (pre) {Tiền xử lý\\resize rộng 800, xám,\\CLAHE, lọc nhiễu};
  \node[box, right=of pre] (edge) {Định vị theo cạnh\\Sobel-x, Otsu, close/open,\\contour, lọc, NMS};
  \node[box, right=of edge] (svm) {HOG + SVM\\(chỉ \code{edge\_svm})\\xếp lại top-5};
  \node[box, below=of svm] (top1) {Chọn ứng viên top-1\\(bbox + quad)};
  \node[box, left=of top1] (quad) {Ước lượng quad\\approxPolyDP trên\\mặt nạ xanh};
  \node[box, left=of quad] (warp) {Homography\\về ảnh $220\times 70$};
  \node[box, below=of warp] (binz) {Nhị phân hóa\\(Otsu hoặc adaptive)};
  \node[box, right=of binz] (seg) {Tách ký tự\\(CCA hoặc projection)};
  \draw[arr] (in) -- (pre);
  \draw[arr] (pre) -- (edge);
  \draw[arr] (edge) -- (svm);
  \draw[arr] (svm) -- (top1);
  \draw[arr] (top1) -- (quad);
  \draw[arr] (quad) -- (warp);
  \draw[arr] (warp) -- (binz);
  \draw[arr] (binz) -- (seg);
\end{tikzpicture}
\caption[Pipeline]{Pipeline. Phương pháp \code{edge} đi thẳng từ khối định vị sang chọn top-1, bỏ qua khối HOG + SVM.}
\label{fig:pipeline}
\end{figure}

\section{Tiền xử lý}
\label{sec:preprocess}

Ảnh gốc được đưa về chiều rộng 800 pixel, giữ tỉ lệ (nội suy diện tích khi thu nhỏ, tuyến tính khi phóng to), rồi đổi sang ảnh xám. Sau đó áp dụng CLAHE \cite{clahe} (giới hạn tương phản 2.0, lưới $8\times 8$) để cân bằng sáng cục bộ, và cuối cùng là một bộ khử nhiễu. Bốn bộ lọc được so sánh ở thí nghiệm E2:
\begin{itemize}
  \item \code{none}: không lọc;
  \item \code{gaussian}: Gauss, kernel $5\times 5$;
  \item \code{median}: trung vị, kernel $5\times 5$;
  \item \code{bilateral}: song phương \cite{bilateral}, đường kính 11, $\sigma_{\text{color}}=17$, $\sigma_{\text{space}}=17$. Bộ lọc này làm mịn vùng phẳng nhưng giữ cạnh có chênh lệch độ sáng lớn.
\end{itemize}
Mọi tham số của bước định vị phía sau (kernel hình thái, ngưỡng diện tích) tính theo ảnh đã resize về 800 pixel nên không phụ thuộc kích thước ảnh gốc.

\section{Phân tích hệ màu}
\label{sec:color}

\textbf{Câu hỏi.} Kênh màu nào tách vùng biển khỏi phần còn lại của ảnh tốt nhất? Câu trả lời cho biết có nên định vị trên ảnh xám (như \code{localize\_edge} hiện tại) hay dùng thêm thông tin màu.

\textbf{Cách đo.} Với mỗi ảnh tune, tính histogram chuẩn hóa 64 bin của từng kênh trong bbox biển thật ($p$) và trên phần ảnh còn lại ($q$). Độ tách là khoảng cách Bhattacharyya
\[
D(p,q)=\sqrt{1-\sum_{i}\sqrt{p_i\,q_i}},
\]
bằng 0 khi hai phân bố trùng nhau và bằng 1 khi không chồng lấn. Bảy kênh được so sánh: ảnh xám, ba kênh HSV (H lấy trong $[0,180)$ theo quy ước OpenCV) và ba kênh Lab. Bảng~\ref{tab:e1} cho trung bình $D$ theo nhóm; E1 luôn chạy trên tune vì chỉ là phân tích dữ liệu, không chọn tham số cho pipeline.

\section{Định vị dựa trên cạnh (\code{edge})}
\label{sec:edge}

Ý tưởng: vùng biển số chứa nhiều nét đứng dày đặc (cạnh các ký tự), nên đạo hàm theo trục $x$ lớn tập trung ở một vùng chữ nhật nằm ngang. Các bước:
\begin{enumerate}
  \item \textbf{Cạnh.} Tính $|G_x|$ bằng Sobel theo $x$ (kernel 3), chuẩn hóa về $[0,255]$ rồi nhị phân hóa bằng Otsu \cite{otsu}.
  \item \textbf{Hình thái học.} Phép \emph{close} với kernel chữ nhật nằm ngang ($21\times 5$, rộng $\times$ cao) để nối các nét ký tự thành một khối, sau đó phép \emph{open} $3\times 3$ để bỏ nhiễu nhỏ.
  \item \textbf{Đa tỉ lệ.} Lặp lại bước close/open với hai kernel phụ $13\times 3$ và $33\times 7$. Biển bị vỡ thành mảnh ở kernel này có thể liền khối ở kernel khác. Ở mỗi kernel, ứng viên được lọc, chấm điểm heuristic và cắt xuống $k=5$ ứng viên cao nhất \emph{trước} khi gộp; sau đó ứng viên của ba kernel được gộp chung, NMS trên bbox (ngưỡng IoU 0.5, giữ điểm cao hơn) và giữ lại $k=5$ ứng viên cao nhất.
  \item \textbf{Ứng viên.} Với mỗi contour ngoài của mặt nạ, lấy hình chữ nhật xoay nhỏ nhất (\code{minAreaRect}); bbox ngang là hình bao của nó, quad là bốn đỉnh sắp theo thứ tự trên-trái, trên-phải, dưới-phải, dưới-trái.
  \item {\sloppy \textbf{Lọc theo đặc trưng hình học.} Với mỗi ứng viên tính: tỉ lệ cạnh dài/ngắn (\code{aspect}), diện tích hình chữ nhật so với diện tích ảnh (\code{area\_ratio}), tỉ lệ pixel mặt nạ trong hình chữ nhật (\code{fill\_ratio}), tỉ lệ pixel cạnh nhị phân trong hình chữ nhật (\code{edge\_density}) và góc nghiêng. Giữ ứng viên thỏa các ngưỡng ở Bảng~\ref{tab:edgeparams}.\par}
  \item \textbf{Điểm heuristic.} Điểm của ứng viên là
  \begin{equation}
    s = 0.4\,d + 0.3\,f + 0.3\exp\!\left(-\frac{\ln^2(a/3.14)}{2\cdot 0.3^2}\right),
  \end{equation}
  trong đó $d$ là \code{edge\_density}, $f$ là \code{fill\_ratio}, $a$ là \code{aspect}; số hạng cuối bằng 1 khi $a\approx 440/140\approx 3.14$, đúng tỉ lệ biển xanh. Giữ tối đa $k=5$ ứng viên điểm cao nhất, quy đổi tọa độ về ảnh gốc bằng cách chia cho hệ số resize.
\end{enumerate}

\begin{table}[htbp]
\centering
\caption{Tham số của bước định vị dựa trên cạnh (cấu hình chốt, tính trên ảnh rộng 800 px).}
\label{tab:edgeparams}
\small
\rowcolors{2}{white}{tablerowalt}
\arrayrulecolor{tableheader}
\begin{tabularx}{\textwidth}{lY}
\hline
\rowcolor{tableheader} \tblhd{Tham số} & \tblhd{Giá trị} \\
\hline
Sobel & theo $x$, kernel 3 \\
close chính / phụ (rộng $\times$ cao) & $21\times 5$ / $13\times 3$, $33\times 7$ \\
open & $3\times 3$ \\
NMS IoU (bbox) & 0.5 \\
\code{aspect} & $[1.5,\ 7.0]$ \\
\code{area\_ratio} & $[0.008,\ 0.16]$ \\
\code{fill\_ratio} tối thiểu & 0.4 \\
\code{edge\_density} tối thiểu & 0.12 \\
Góc nghiêng tối đa & $30^\circ$ \\
Số ứng viên giữ lại $k$ & 5 \\
\hline
\end{tabularx}
\end{table}

Thuật toán~\ref{alg:edge} tóm tắt toàn bộ bước định vị theo cạnh.

\begin{algorithm}[htbp]
\caption{Định vị biển dựa trên cạnh (\code{edge})}
\label{alg:edge}
\begin{algorithmic}[1]
\Require Ảnh xám đã tiền xử lý $I$ (rộng 800 px), số ứng viên giữ lại $k=5$
\Ensure Tối đa $k$ ứng viên (bbox, quad, điểm), sắp theo điểm giảm dần
\State $G \gets |\textsc{Sobel}_x(I)|$ chuẩn hóa về $[0,255]$;\ $M_0 \gets \textsc{Otsu}(G)$
\State $C \gets \emptyset$
\ForAll{kernel close $K \in \{21\times 5,\ 13\times 3,\ 33\times 7\}$}
    \State $M \gets \textsc{Open}_{3\times 3}\big(\textsc{Close}_K(M_0)\big)$
    \State $C_K \gets \emptyset$
    \ForAll{contour ngoài của $M$}
        \State $r \gets \textsc{MinAreaRect}$;\ tính \code{aspect}, \code{area\_ratio}, \code{fill\_ratio}, \code{edge\_density}, góc
        \If{$r$ thỏa các ngưỡng ở Bảng~\ref{tab:edgeparams}}
            \State thêm $r$ vào $C_K$ với điểm heuristic $s$
        \EndIf
    \EndFor
    \State $C \gets C \cup \textsc{TopK}(C_K, k)$
\EndFor
\State $C \gets \textsc{NMS}(C,\ \text{IoU}=0.5)$;\ $C \gets \textsc{TopK}(C, k)$
\State quy đổi tọa độ về ảnh gốc (chia cho hệ số resize);\ \Return $C$
\end{algorithmic}
\end{algorithm}

Các ngưỡng này được chọn bằng cách quét trên tập tune (mục~\ref{sec:tuning}). Riêng cận dưới của \code{area\_ratio} bám theo phân bố diện tích bbox thật trên tune (phân vị 5\% là 0.0116, lớn nhất 0.153), nhằm loại logo và chữ nhỏ.

\section{Chấm điểm lại ứng viên bằng HOG và SVM (\code{edge\_svm})}
\label{sec:hogsvm}

Điểm heuristic chỉ gồm mật độ cạnh, độ lấp đầy và tỉ lệ cạnh nên dễ xếp nhầm chữ hãng xe hay logo lên trên biển thật. Bước này thay điểm đó bằng điểm của bộ phân loại học từ ảnh:
\begin{itemize}
  \item \textbf{Đặc trưng.} Cắt bbox của ứng viên trên ảnh xám, resize về $128\times 32$ (rộng $\times$ cao), trích HOG \cite{hog}: 9 hướng, ô $8\times 8$ pixel, khối $2\times 2$ ô, chuẩn hóa L2-Hys, vector $D=1620$ chiều.
  \item \textbf{Mẫu dương.} Crop bbox thật của 198 ảnh tune.
  \item \textbf{Mẫu âm.} Với mỗi mẫu dương lấy 3 mẫu âm: ưu tiên ứng viên sai của \code{localize\_edge} có IoU với bbox thật dưới 0.3, nếu thiếu thì bù bằng vùng ngẫu nhiên (rộng 8--45\% chiều rộng ảnh, tỉ lệ cạnh trong \code{aspect}, IoU dưới 0.3). Tổng cộng 198 dương và 594 âm.
  \item \textbf{Bộ phân loại.} SVM tuyến tính \cite{svm} (\code{LinearSVC}, $C=1$, seed 42). Điểm của ứng viên là giá trị hàm quyết định; các ứng viên được sắp giảm dần theo điểm này. Không dùng ngưỡng điểm để loại ứng viên, nhằm tính được Recall@$k$.
  \item \textbf{Kiểm chứng chéo.} 5 fold phân tầng trên tập mẫu của tune để ước lượng chất lượng phân loại patch.
\end{itemize}
Thuật toán~\ref{alg:svm} tóm tắt huấn luyện và chấm điểm lại.

\begin{algorithm}[htbp]
\caption{Huấn luyện SVM và xếp lại ứng viên (\code{edge\_svm})}
\label{alg:svm}
\begin{algorithmic}[1]
\Require Ảnh tune (198 ảnh) kèm bbox thật; tối đa $k=5$ ứng viên của \code{localize\_edge} cho ảnh cần xử lý
\Ensure Danh sách ứng viên sắp theo điểm SVM giảm dần
\Procedure{Train}{}
    \ForAll{ảnh tune}
        \State mẫu dương: crop bbox thật, resize $128\times 32$, trích HOG ($D=1620$)
        \State 3 mẫu âm: ứng viên sai có IoU $<0.3$, thiếu thì bù vùng ngẫu nhiên có IoU $<0.3$
    \EndFor
    \State huấn luyện \code{LinearSVC} ($C=1$, seed 42) trên 198 dương và 594 âm
\EndProcedure
\Procedure{Rerank}{ứng viên}
    \ForAll{ứng viên}
        \State điểm $\gets$ giá trị hàm quyết định SVM trên HOG của crop bbox ($128\times 32$)
    \EndFor
    \State \Return ứng viên sắp theo điểm giảm dần (không loại ứng viên nào)
\EndProcedure
\end{algorithmic}
\end{algorithm}

SVM chỉ xếp lại tối đa 5 ứng viên mà bước cạnh đã sinh ra, nên nó không thể cứu ảnh mà bước cạnh không sinh ứng viên đúng. Mô hình chỉ huấn luyện trên tune; mã nguồn từ chối dữ liệu không thuộc tune.

\section{Nắn phối cảnh}
\label{sec:rectify}

Ứng viên top-1 cho một bbox ngang và một quad. Biển nghiêng hoặc chụp xiên có quad không phải hình chữ nhật ngang, nên cắt bbox giữ lại nền và làm ký tự méo. Phép nắn đưa biển về ảnh chuẩn $220\times 70$ (tỉ lệ $\approx 3.14$, khớp biển xanh $440\times 140$ mm). Có hai cách lấy quad:
\begin{itemize}
  \item \code{minarearect}: dùng luôn bốn đỉnh \code{minAreaRect} của ứng viên;
  \item \code{approxpoly} (mặc định): tinh chỉnh bốn góc từ màu biển. Trong bbox nới thêm 15\%, tạo mặt nạ màu xanh dương trong HSV ($H\in[95,135]$, $S\ge 70$, $V\ge 40$), đóng và mở hình thái, giới hạn mặt nạ trong quad ứng viên (tránh rò sang nền xanh), lấy contour lớn nhất chứa tâm bbox, bao lồi, rồi \code{approxPolyDP} với $\varepsilon$ tăng từ 1\% đến 15\% chu vi (bước 0.5\%) tới khi còn đúng 4 đỉnh (dừng sớm nếu số đỉnh đã tụt xuống dưới 4). Nếu không ra 4 đỉnh, quad không lồi hoặc diện tích ngoài khoảng $30\%$--$300\%$ diện tích bbox thì lùi về quad \code{minAreaRect}.
\end{itemize}
Với bốn góc nguồn $\mathbf{p}_i$ (sắp theo thứ tự trên-trái, trên-phải, dưới-phải, dưới-trái) và bốn góc đích $(0,0)$, $(219,0)$, $(219,69)$, $(0,69)$, ma trận homography $H\in\mathbb{R}^{3\times 3}$ thỏa $\tilde{\mathbf{x}}'\sim H\tilde{\mathbf{x}}$ (tọa độ thuần nhất) được giải từ bốn cặp điểm, rồi áp dụng nội suy tuyến tính, viền lặp pixel biên. Phương án đối chứng là cắt bbox ngang rồi resize về $220\times 70$ (\code{bbox}), không bù được độ nghiêng.

\section{Nhị phân hóa: Otsu và adaptive}
\label{sec:binarize}

Đầu vào của bước nhị phân hóa là ảnh biển đã nắn về kích thước $220\times 70$ (mục~\ref{sec:rectify}), chuyển sang ảnh xám. Hàm \code{binarize\_plate} nhận ảnh biển xám này và trả ảnh nhị phân với \emph{ký tự trắng trên nền đen}, theo một trong hai cách.

\textbf{Otsu~\cite{otsu}.} Chọn một ngưỡng toàn cục $t$ cho cả biển, sao cho phương sai giữa hai lớp (pixel $\le t$ và $>t$) lớn nhất:
\[
t^{*}=\arg\max_{t}\;\omega_0(t)\,\omega_1(t)\,\big(\mu_0(t)-\mu_1(t)\big)^2,
\]
trong đó $\omega_k$ là tỉ lệ pixel và $\mu_k$ là độ sáng trung bình của lớp $k$. Cách này hợp khi histogram của biển có hai đỉnh rõ (nền và chữ) và độ sáng đều trên toàn biển.

\textbf{Adaptive.} Mỗi pixel có ngưỡng riêng $T(x,y)=m_G(x,y)-C$, với $m_G$ là trung bình có trọng số Gauss trong cửa sổ $25\times 25$ quanh pixel và $C=7$ (tham số \code{binarize.adaptive} trong config). Cách này chịu được biển sáng không đều (bóng đổ, chói một bên). Bản cài đặt đầu tiên dùng \code{THRESH\_BINARY} trực tiếp: với biển xanh chữ trắng, vùng nền phẳng thành trắng còn ký tự chỉ còn viền rỗng. Bản sửa dùng Otsu để đoán cực tính (nếu chữ sáng hơn nền thì đảo ảnh xám trước) rồi áp \code{THRESH\_BINARY\_INV}, nên nền phẳng luôn thành đen; một test đơn vị kiểm tra điều này.

Sau đó \code{ensure\_text\_white} đảo ảnh nếu hơn một nửa số pixel là trắng (phòng trường hợp đoán sai cực tính), và \code{clean\_binary} mở hình thái học bằng kernel $2\times 2$ để xóa chấm nhiễu nhỏ.

\section{Tách ký tự: CCA và projection}
\label{sec:segment}

Cả hai bộ tách nhận ảnh nhị phân (chữ trắng) và trả danh sách hộp ký tự sắp từ trái sang phải. Các ngưỡng hình học tính theo tỉ lệ kích thước biển nên không phụ thuộc độ phân giải; các tham số thí nghiệm (khoảng chiều cao ký tự $[0.4,0.9]$ chiều cao biển, tỉ lệ rộng/cao $[0.10,1.0]$, diện tích tối thiểu 30 px) lấy từ config.

\textbf{CCA (connected component analysis).} Gán nhãn thành phần liên thông 8 hướng, rồi:
{\sloppy
\begin{enumerate}
  \item bỏ viền và nhiễu: thành phần rộng hơn 30\% biển mà rỗng (tỉ lệ lấp đầy dưới 0.35) là viền biển; thành phần có tâm theo chiều dọc ngoài khoảng $[0.25,0.75]$ chiều cao là ốc vít hoặc mép. Khối rộng nhưng đặc (nhiều ký tự dính nhau) được \emph{giữ lại} để cắt ở bước 4, thay vì bị xóa như viền;
  \item gộp các mảnh chồng nhau theo chiều ngang (chồng $\ge 50\%$), vì ký tự Hán thường vỡ thành nhiều nét xếp chồng; sau đó mới bỏ mảnh dẹt đứng riêng (cao dưới 20\% biển);
  \item gộp mảnh hẹp (rộng/cao $<0.15$) với mảnh kề nếu khe đủ nhỏ, để không mất các nét dọc rời của ký tự Hán;
  \item lọc theo chiều cao, tỉ lệ và diện tích, rồi \code{merge\_split\_boxes}: gọi $w_{\mathrm{med}}$ là trung vị chiều rộng hộp và $g$ là trung vị khe giữa các hộp kề; hai hộp hẹp ($<0.75\,w_{\mathrm{med}}$) gần nhau được gộp nếu hộp gộp không quá $1.4\,w_{\mathrm{med}}$; hộp rộng $w\ge 1.7\,w_{\mathrm{med}}$ được cắt đều thành $\operatorname{round}\big((w+g)/(w_{\mathrm{med}}+g)\big)$ phần (tối đa 3).
\end{enumerate}
\par}

\textbf{Projection.} Sau khi bỏ viền và ốc vít theo cùng luật, \code{split\_lines} chiếu theo hàng để tách dòng (CCPD chỉ có biển một dòng; trường hợp hai dòng được kiểm bằng test giả). Trong mỗi dòng, profile cột là số pixel trắng của từng cột; cột có ít nhất 10\% chiều cao dòng là trắng được coi là cột chữ. Các dải cột chữ cách nhau dưới 2 px được nối, dải hẹp được nối với dải kề (ký tự Hán vỡ), dải dưới 4 px bị bỏ. Dải rộng $\ge 1.7\,w_{\mathrm{med}}$ được cắt tại cột có ít pixel trắng nhất (``thung lũng'') trong khoảng $\pm 0.3\,w_{\mathrm{med}}$ quanh vị trí cắt dự kiến, thay vì cắt đều như CCA. Cuối cùng mỗi hộp được co theo chiều dọc về phần có chữ và lọc như CCA.
